\section{Transport de l'audio}
\label{sec:transport}

Après avoir détaillé le traitement du signal par le \gls{sink}, cette section décrit son acheminement depuis la carte \gls{sd}. Deux exigences guident cette conception. Il s'agit de garantir un flux continu adapté à chaque fichier sans interruption audible et de maintenir l'indépendance vis-à-vis de la sortie active pour faciliter la commutation jack/Bluetooth.

Le transport s'articule autour de quatre maillons séquentiels. Le service de stockage gère l'index des fichiers sur la carte \gls{sd}. Les décodeurs convertissent ces fichiers en un flux \gls{pcm} uniforme. Le pipeline achemine ce flux vers le \gls{sink} actif sous le contrôle de la tâche audio. Enfin l'orchestration gère la sélection des pistes et leur enchaînement.

\subsection{Source de données et indexation}

Le service de stockage possède le montage de la carte SD, réalisé au travers de la couche \gls{fatfs} d'ESP-IDF~\cite{ESP32_IDF_guide}. Lorsqu'une carte est présente, il en parcourt récursivement l'arborescence, retient les fichiers d'extension \texttt{.mp3} ou \texttt{.wav} sans distinction de casse, et trie le résultat par ordre alphabétique. Les entrées dont le nom commence par un point sont écartées, ce qui neutralise les fichiers cachés et les fichiers de ressources que certains systèmes d'exploitation ajoutent.

Ce parcours récursif ne fixe aucune limite de profondeur explicite, la seule contrainte réelle étant la taille de la pile d'exécution à chaque appel. Le chemin complet est limité à 320 caractères et toute entrée dépassant cette longueur est ignorée sans interrompre le balayage.

L'index réside en \gls{psram} pour préserver la mémoire interne (section~\ref{sec:archi}). Il occupe un bloc contigu de 512 entrées de 320 octets. Un dépassement provoque un échec explicite plutôt qu'une troncature silencieuse, signalant ainsi clairement une carte trop remplie. Une carte vide mais lisible constitue un succès avec zéro piste, distinct d'une erreur de lecture.

Le démontage de la carte ne libère pas le bus mais retire uniquement le périphérique associé. Bien que l'écran et la carte disposent d'hôtes \gls{spi} dédiés, les contrôleurs SPI partagent seulement deux canaux DMA pour trois hôtes~\cite{ESP32_Technical_Reference_Manual}. Libérer le bus suspendrait un transfert d'affichage en cours et bloquerait l'interface, un comportement observé lors des tests. Le démontage partiel écarte ce risque tout en suffisant à la gestion des insertions et retraits.

La tâche de maintenance gère à chaud les changements de carte en surveillant la ligne de détection. Une transition n'est retenue qu'après deux relevés concordants pour ignorer les rebonds de contact. Un retrait stoppe la lecture puis démonte le système de fichiers, tandis qu'une insertion déclenche une réindexation confiée à une tâche éphémère.

\subsection{Décodage et représentation interne}

Le décodage s'appuie sur une façade unique offrant une interface générique pour l'ouverture, la description, la lecture et la fermeture des fichiers. Cette couche oriente chaque fichier vers un décodeur \gls{mp3} ou \gls{wav} selon sa signature réelle, privilégiant le contenu de l'entête à l'extension pour éviter les erreurs de nommage. Le décodage MP3 utilise la bibliothèque libhelix fonctionnant en virgule fixe~\cite{helix_mp3_decoder}, éliminant ainsi le recours coûteux à l'unité de calcul flottant.

Cette abstraction se justifie par la prise en charge simultanée du MP3 et du WAV exigée par le cahier des charges. Un lecteur limité au seul MP3 n'aurait nécessité qu'un module dédié, mais la présence de deux formats rend la façade indispensable.

Les décodeurs produisent une représentation interne unique sur 32 bits justifiés à gauche et entrelacés (section~\ref{sec:audio_filaire}). Le MP3 et le WAV 16 bits subissent un décalage vers les poids forts, tandis que le WAV 24 bits positionne ses trois octets utiles sur les bits de poids fort. La lecture s'effectue donc octet par octet pour gérer les échantillons de 24 bits, reflétant symétriquement l'écriture vers le sink.

La durée des pistes est exacte pour le WAV grâce à son en-tête. Pour le MP3, elle dépend de l'en-tête de débit variable si présent. En l'absence de cette information, la durée est déclarée inconnue plutôt qu'estimée, évitant ainsi toute erreur liée à une hypothèse de débit constant sur un fichier à débit variable.

\subsection{Pipeline audio}

Le pipeline est porté par la tâche audio, seule productrice du sink actif. Conformément au modèle d'exécution de la section~\ref{sec:archi}, elle demeure bloquée au repos en attente d'une commande. Son interface n'est pas bloquante pour l'appelant, chaque ordre de lecture, d'arrêt, de pause, de reprise ou de bascule de sortie étant déposé dans une file à quatre entrées que la boucle consomme en tête de cycle~\cite{freertos_kernel_doc}.

Le décodage produit les échantillons plus vite qu'ils ne sont joués, c'est donc la sortie qui impose le rythme. Le pipeline lui remet un bloc d'échantillons à chaque tour de boucle. La sortie filaire l'absorbe en entier, mais la sortie Bluetooth, dont le tampon interne se vide à la cadence de la liaison radio, peut n'en prendre qu'une partie lorsqu'il est momentanément plein. Le pipeline conserve alors la portion restante et la présente de nouveau au tour suivant, si bien qu'aucun échantillon n'est perdu et que le débit s'aligne de lui-même sur celui de la sortie. Comme les commandes en attente sont traitées au début de chaque tour de boucle, un ordre d'arrêt ou de pause reste pris en compte sans délai, même si la sortie est saturée à cet instant.

La position de lecture est mise à jour en continu par la tâche audio et lue par la tâche d'interface pour afficher la progression. Comme cette position tient dans un entier de 32~bits, le processeur l'écrit et la lit en une seule opération indivisible. La tâche d'interface ne peut donc jamais surprendre une valeur à moitié mise à jour, elle lit toujours soit l'ancienne, soit la nouvelle, ce qui rend inutile tout verrou de protection. La position est remise à zéro à l'arrêt et figée en pause, une valeur nulle signalant alors qu'aucune piste ne joue.

\subsection{Orchestration de la lecture}

Un orchestrateur sans tâche dédiée relie l'interface, la logique de liste de lecture et le stockage. Il maintient l'état du lecteur (arrêté, lecture ou pause) et soumet des commandes au pipeline sans accéder directement au \gls{sink}. La pause stoppe le \gls{sink} tout en conservant le décodeur ouvert et la position courante, permettant une reprise instantanée.

L'enchaînement des pistes est automatique. La fin d'une piste déclenche un rappel dans le contexte de la tâche audio, permettant à l'orchestrateur de sélectionner la suivante selon le mode de répétition. Chaque transition rouvre le décodeur pour la nouvelle piste. De plus, aucun envoi d'échantillon n'est autorisée avant l'acquittement du volume, une vérification systématique détaillée à la section~\ref{sec:soft_bluetooth}.

La liste de lecture conserve un état minimal composé de la position courante, du mode de répétition et d'un ordre aléatoire précalculé. Elle constitue une vue dynamique des pistes à venir plutôt qu'une structure dupliquée. Seules les pistes futures peuvent être réordonnées ou retirées, les pistes jouées ou en cours restant figées. Un retrait n'efface pas le fichier de la carte SD et la piste réapparaît lors d'une prochaine indexation. Cette réindexation met à jour la liste et corrige la position si l'index a évolué.

La bascule de sortie en cours de lecture préserve la position et la progression. La sortie active continue de jouer jusqu'à ce que la nouvelle soit prête, le pipeline reprenant alors sur le \gls{sink} choisi. Le cas spécifique du passage vers le Bluetooth est traité à la section~\ref{sec:soft_bluetooth}.